\section{数据集重建与标注规则}

\subsection{筛选规则}

本文使用的子集由 \path{src/prepare_mswc_initial_gbdz_subset.py} 生成，规则很简单：
\begin{quote}
保留一个单词，当且仅当它的首字母属于 \texttt{g/b/d/z}，并且 CMUdict / ARPABET 中首个非元音音素与该字母严格对应为 \texttt{G/B/D/Z}。
\end{quote}

这条规则的作用，是把“拼写里有目标字母”收紧成“这个目标辅音确实出现在起始位置”。这样一来，数据不再混入发音位置不对的样本，也不再把字母出现位置和实际发音边界混在一起。

\subsection{数据规模}

\begin{table}[H]
    \centering
    \caption{精确首浊辅音子集的规模统计}
    \label{tab:dataset}
    \begin{tabular}{lcccc}
        \toprule
        标签 & 词数 & train & dev & test \\
        \midrule
        /g/ & 15 & 600 & 75 & 75 \\
        /b/ & 15 & 600 & 75 & 75 \\
        /d/ & 15 & 600 & 75 & 75 \\
        /z/ & 15 & 600 & 75 & 75 \\
        \midrule
        合计 & 60 & 2400 & 300 & 300 \\
        \bottomrule
    \end{tabular}
\end{table}

这个子集是完全平衡的，测试集每类 75 条，总计 300 条。相比旧版更松的词级设置，它更接近“首浊辅音检测”这个题目本身。
